\section{本章小结}
\label{sec:ls-dual-spaces-summary}

本章回答了章首的三个问题：线性读取组成对偶空间，读取经对偶映射在两端之间传递，读取的大小由对偶范数度量。

对偶空间由线性泛函组成，维数与原空间相同；对偶基就是坐标读取，用对偶基表示一个读取时不必解方程。对偶映射把输出端的读取传回输入端，在对偶基下记录为转置，由此又得到行秩等于列秩。这些结论只用线性运算，与范数和内积无关。

有了范数，读取就有了大小。对偶范数是读取在单位球上的最大绝对值，也就是误差受限时读取的最大变化；$\ell_p$范数的对偶范数是$\ell_q$范数，$1/p+1/q=1$，最坏方向由Hölder不等式的取等条件给出。讨论“输入误差最多引起多大输出误差”时，必须先说明误差用哪种范数度量。

有了内积，读取又可以写成向量。Riesz表示定理说明，有限维内积空间上的每个线性泛函都是与唯一一个向量的内积，内积诱导范数下对偶范数就是这个向量的长度。把对偶映射两端的读取都换成向量，就得到伴随：在标准正交基下伴随就是转置，在一般基下还要计入两端的Gram矩阵。自伴、正交与正规算子都由伴随定义，第~\ref{chap:eigen-spectral-decomposition}章的谱定理将在这类算子上展开。
